\chapter{Introduction}
\label{chap:introduction}

\section{Motivation}

Fibrous particles occur in many natural and industrial suspensions, including
pulp fibres, biological filaments, textile fibres, and fibrous microplastics.
This thesis is motivated primarily by the last of these examples. Synthetic
microfibres are now commonly reported in marine water and sediments, and their
large aspect ratios distinguish them from the spherical or compact particles
often used as model contaminants \cite{gago_synthetic_2018}. Their transport is
therefore not determined only by size and density. Orientation, bending
stiffness, hydrodynamic torque, and deformation all influence how a fibre
settles, reorients, samples the surrounding flow, and interacts with nearby
fibres.

For ocean transport problems, this matters because the relevant particles are
small but not necessarily dynamically simple. A fibre may be dilute enough that
direct contact with other fibres is rare, while still being part of a
suspension large enough that one-particle simulations are not representative.
Likewise, a single settling velocity is not enough to describe the motion when
the fibre can rotate, bend, or align with a local shear. These effects become
especially important when the particle Reynolds number is not asymptotically
small. In that regime, fluid inertia can alter the drag, torque, wake
structure, and collective settling behaviour, so a purely Stokesian description
may miss the mechanisms controlling vertical transport.

The computational difficulty is that the physically relevant model contains
several coupled pieces: an incompressible fluid, deformable inextensible
fibres, no-slip coupling between the two, and a geometry that changes in time.
A fully body-fitted description of every fibre would be expensive even for a
single object, and becomes impractical when the target application is a dilute
many-fibre suspension. Computational efficiency is therefore a central
motivation for the model developed here. The goal is not to resolve every
microscopic detail of a fibre surface, but to retain the leading mechanics
needed for settling and suspension dynamics while keeping the method scalable
enough to simulate many moving fibres.

\section{Background}

The dynamics of flexible fibres in viscous flows have been studied through a
combination of asymptotic theory, numerical simulation, and laboratory
experiment. At low Reynolds number, the slender geometry of a fibre motivates
centerline descriptions and slender-body approximations in which the
three-dimensional hydrodynamic interaction is reduced to forces distributed
along a curve. This has led to a substantial literature on the deformation,
buckling, rotation, and suspension behaviour of flexible filaments in viscous
flows; a broad review is given by \textcite{du_roure_dynamics_2019}. These
results show that even in the absence of inertia, flexibility introduces
qualitatively new behaviour compared with rigid rods.

Settling problems make this coupling particularly clear. In the Stokes regime,
\textcite{marchetti_deformation_2018} combined experiments, numerical
simulations, and scaling arguments to identify distinct deformation regimes for
a single flexible fibre settling in a quiescent viscous fluid. In weakly
deformed and strongly deformed limits the drag remains proportional to the
settling speed, while an intermediate elastic reconfiguration regime changes
the effective drag law. This is an example of the general mechanism that
motivates the present work: fibre shape is part of the transport problem, not
merely an output of it.

The literature is less complete once fluid inertia and many interacting fibres
are considered together. There are numerical studies of flexible fibre
suspensions at finite inertia, including the inertial settling simulations of
\textcite{banaei_inertial_2020}, but this regime remains much less developed
than the single-fibre and low-Reynolds-number cases. The gap is partly physical
and partly computational. Inertial simulations require a Navier--Stokes solver
capable of resolving wakes and unsteady hydrodynamic forces, while
many-particle simulations require a coupling method whose cost does not grow
prohibitively with the number of fibres. The modelling choices in this thesis
are aimed at this intersection: finite-Reynolds-number flow, flexible
inextensible fibres, and dilute suspensions large enough that efficiency is a
limiting design constraint.

\section{Problem Statement}

The physical system considered in this thesis is a dilute suspension of
slender, flexible fibres immersed in an incompressible Newtonian fluid. The
fluid is described by the Navier--Stokes equations, while each fibre is
represented by a Lagrangian centerline. The structural model treats the fibre
as an Euler--Bernoulli beam with large rotations and an inextensibility
constraint. The fluid and fibre are coupled by a no-slip condition: the
velocity of the fibre centerline must agree with the local fluid velocity, and
the force needed to enforce this condition is applied back to the fluid.

This formulation involves two constraint mechanisms. The fluid velocity must
remain divergence-free, which introduces pressure as a constraint force for
incompressibility. The fibre centerline must remain inextensible, which
introduces internal constraint forces along the beam. In addition, the immersed
boundary coupling imposes a kinematic constraint between the Eulerian fluid
velocity and the Lagrangian fibre velocity. A useful numerical method must
therefore handle constraints at both the fluid and structural levels without
making each time step too expensive.

The central modelling trade-off is between fidelity and cost. A surface- or
volume-resolved fibre model would represent the no-slip boundary more sharply,
but would also require substantially more degrees of freedom per fibre and a
more complicated treatment of moving geometry. The centerline immersed-boundary
model used here instead regularizes the fibre-fluid interaction over a small
Eulerian stencil. This introduces a controlled loss of geometric sharpness, but
it makes the method compatible with fixed Cartesian grids, adaptive mesh
refinement, and matrix-free force-transfer operations. These features are
important for the long-term target of many-fibre simulations.

\section{Approach}

The thesis develops the model in modular pieces, which reflects the overall
partitioned nature of the scheme and the implementation itself. The fibre
mechanics are built from large-deflection Euler--Bernoulli beam theory.
Inextensibility is enforced using an alternating-direction method of
multipliers, following the constrained beam formulation of Bourgat, Dumay, and
Glowinski \cite{bourgat_large_1980,glowinski_large_nodate}. The splitting
introduces a tangent variable so that the nonlinear unit-length constraint
becomes a local projection, while the global centerline update remains a linear
beam solve. This structure is attractive computationally because the expensive
matrix factorization can be reused during nonlinear iterations. Alternative
discretization schemes for the beam model are also discussed.

The fluid solver is based on a projection method for the incompressible
Navier--Stokes equations on adaptive Cartesian grids. Tree-structured adaptive
mesh refinement allows the mesh to concentrate resolution near fibres, wakes,
and other localized flow features, while leaving the rest of the domain coarse.
This is another efficiency choice: resolving a fibre suspension uniformly at
the finest required scale would waste most of the mesh away from the particles.

The fluid--fibre coupling is formulated as an immersed-boundary method. A
regularized delta kernel transfers velocity from the Eulerian grid to
Lagrangian markers and spreads marker forces back to the fluid. Direct forcing,
multi-direct forcing, and a marker-space distributed-Lagrange-multiplier
interpretation are used to explain the coupling force as an approximate
constraint force enforcing no-slip. The resulting method avoids body-fitted
remeshing, handles large fibre motion naturally, and keeps the coupling
operations local except for the fluid projection and any optional marker-space
linear solve.

\section{Thesis Organization}

Chapter~\ref{chap:euler-bernoulli} derives the beam model used for the fibre
centerline, beginning with the Euler--Bernoulli assumptions and then moving to
the large-rotation inextensible formulation. The alternating-direction
structural solver and its validation are described in Chapter~\ref{chap:addm}.
Chapter~\ref{chap:navier-stokes} summarizes the adaptive-grid Navier--Stokes
discretization and the projection solver used for the fluid.
Chapter~\ref{chap:immersed-boundary-methods} develops the immersed-boundary
coupling, including direct forcing, multi-direct forcing, and the marker-space
Lagrange-multiplier viewpoint. Finally, Chapter~\ref{chap:coupled-validation}
assembles these pieces for validation of the coupled fluid--fibre system.
